(1) 沿 \(z\) 方向入射的低强度单色平面电磁波照射自由非相对论电子。定义 \(r_e=e^2/(4\pi\varepsilon_0mc^2)\)，辐射观测方向 \(\boldsymbol n=(\sin\theta,0,\cos\theta)\)。从电子加速度 \(\boldsymbol a=-e\boldsymbol E/m\) 和远场偶极辐射的横向投影，推导固定单位入射线偏振 \(\boldsymbol e\) 时的微分散射截面。
(2) 对入射波沿 \(x\)、沿 \(y\) 的两种正交线偏振分别求散射平面内与垂直方向的差分截面。取不偏振入射的平均，推导角分布及总截面 \(\sigma_T\)，需要实际完成球面角积分。
(3) 求散射出射光的线偏振度 \(\Pi(\theta)\)，考察 \(\theta=0,\pi/2,\pi\) 的极限；解释在入射不偏振时为什么90度散射可以完全线偏振。
(4) 设入射能量流密度 \(I_0\)。从散射角分布积分求电子平均吸收的纵向动量通量（辐射压力力）及散射总功率，检查单位，并讨论近相对论频率电子反冲/Compton效应造成的失效条件。
